\documentclass[12pt,a4paper]{article}
\usepackage[utf8]{inputenc}
\usepackage[T1]{fontenc}
\usepackage[spanish]{babel}
\usepackage{amsmath,amssymb}
\usepackage{geometry}
\usepackage{booktabs}
\usepackage{longtable}
\usepackage{graphicx}
\usepackage{array}
\usepackage{pdflscape}
\geometry{margin=2.5cm}

\title{Tarea --- Análisis Comparativo de Conductores Eléctricos\\para Líneas de Transmisión}
\author{Emerson Warhman\\Líneas de Transmisión}
\date{1 de junio de 2026}

\begin{document}
\maketitle

\section{Introducción}

Este informe presenta un análisis comparativo de los principales tipos de
conductores utilizados en líneas de transmisión eléctrica, evaluando sus
parámetros técnicos para el calibre asignado de \textbf{950 MCM}.

La investigación se realizó con base en el catálogo técnico de \textbf{CVG CABELUM}
(Conductores de Aluminio del Caroní, C.A.), empresa estatal venezolana ubicada
en Ciudad Bolívar, única productora nacional de conductores eléctricos de alta
potencia.

\section{Consideraciones de Fabricación y Alcance del Catálogo}

Al revisar exhaustivamente el catálogo de CVG CABELUM, se identificaron las
siguientes consideraciones técnicas importantes:

\begin{enumerate}
    \item \textbf{Gama de conductores desnudos:} CABELUM fabrica comercialmente
        \textbf{AAC} (Aluminio 1350-H19), \textbf{AAAC} (Aleación de Aluminio
        6201-T81), \textbf{ACAR} (Aluminio 1350 reforzado con Aleación 6201) y
        \textbf{ACSR} (Aluminio reforzado con acero galvanizado).

    \item \textbf{AACSR:} \textbf{No está disponible} en las hojas técnicas del
        catálogo de CABELUM. Para efectos académicos, se desglosará su
        comportamiento teórico en base a las físicas de la planta.

    \item \textbf{Normas DIN y BRITISH:} Las tablas técnicas de CABELUM bajo
        norma BRITISH (BS 215 / BS 3242) y DIN (48201) solo llegan hasta
        secciones de 794 mm$^2$ (calibre Tarantula) y 1000 mm$^2$
        respectivamente. \textbf{No existen especificaciones de 950 MCM bajo
        estas normas.} Por lo tanto, la comparativa real de ingeniería se
        efectúa bajo el estándar de diseño de planta: \textbf{ASTM}.

    \item \textbf{Normas aplicables:} ASTM B-231 para AAC, ASTM B-399 para
        AAAC, ASTM B-524 para ACAR y ASTM B-232 para ACSR.
\end{enumerate}

\section{Tipos de Conductores Evaluados}

\subsection{AAC (All Aluminum Conductor) --- ASTM B-231}
Conductor de aluminio puro 1350-H19. Características:
\begin{itemize}
    \item Alta conductividad eléctrica
    \item Bajo peso
    \item Menor resistencia mecánica
    \item Usos: distribución urbana, cortas distancias
\end{itemize}

\subsection{AAAC (All Aluminum Alloy Conductor) --- ASTM B-399}
Conductor de aleación de aluminio 6201-T81. Características:
\begin{itemize}
    \item Mayor resistencia mecánica que AAC
    \item Buena relación resistencia-peso
    \item Resistencia a la corrosión
    \item Usos: líneas de media y larga distancia
\end{itemize}

\subsection{ACAR (Aluminum Conductor Alloy Reinforced) --- ASTM B-524}
Conductor de aluminio con refuerzo de aleación. Características:
\begin{itemize}
    \item Combinación de alambres de aluminio puro y de aleación 6201
    \item Balance entre conductividad y resistencia
    \item Usos: líneas que requieren mayor capacidad mecánica
\end{itemize}

\subsection{AACSR (Aluminum Alloy Conductor Steel Reinforced) --- Teórico}
Conductor de aleación de aluminio con refuerzo de acero. Características:
\begin{itemize}
    \item Alta resistencia mecánica (por el acero)
    \item Mayor peso
    \item \textbf{Nota:} No fabricado por CABELUM. Valores estimados.
    \item Usos: grandes claros, condiciones extremas
\end{itemize}

\subsection{ACSR (Aluminum Conductor Steel Reinforced) --- ASTM B-232}
Conductor de aluminio con refuerzo de acero galvanizado. Características:
\begin{itemize}
    \item El más utilizado en transmisión
    \item Excelente relación resistencia-peso
    \item Máxima resistencia mecánica de la línea
    \item Usos: líneas de transmisión de alta tensión
\end{itemize}

\section{Tabla Comparativa --- Calibre 950 MCM}

Evaluado bajo los parámetros disponibles en el catálogo de CABELUM para el
calibre asignado:

\begin{landscape}
\begin{center}
\begin{tabular}{p{4cm}c|c|c|c}
\toprule
\textbf{Parámetro} & \textbf{AAC} & \textbf{AAAC} & \textbf{ACAR} & \textbf{ACSR} \\
& (Goldenrod) & (Greeley) & (950 kcmil 18/19) & (Rail) \\
\midrule
Norma de fabricación & ASTM B-231 & ASTM B-399 & ASTM B-524 & ASTM B-232 y B-498 \\
\midrule
Calibre comercial & 954 kcmil & 927.2 kcmil & 950.0 kcmil & 954.0 kcmil \\
\midrule
Sección (mm$^2$) & 483.4 & 470.0 & 481.0 & 483.85 \\
& & & & (Al: 450.26 / Acero: 33.59) \\
\midrule
Construcción (hilos) & 61 Al & 37 Aleación & 18 Al / 19 Aleac. & 45 Al / 7 Acero \\
\midrule
Diámetro (mm) & 28.62 & 28.14 & 28.48 & 29.61 \\
\midrule
Peso neto (kg/km) & 1331.0 & 1289.0 & 1321.0 & 1598 \\
& & & & (Al: 1337 / Acero: 261) \\
\midrule
Carga de rotura (kgf) & 7648 & 13767 & 10810 & 11749 \\
\midrule
Resistencia DC 20$^\circ$C ($\Omega$/km) & 0.05964 & 0.07133 & 0.06490 & 0.05988 \\
\midrule
Ampacidad (A) & 985 & 905 & 1044 & 990 \\
\midrule
Long. máx. despacho (m) & 2470 & 2550 & 2490 & 2110 \\
\midrule
Peso neto carrete (kg) & 3288 & 3282 & 3289 & 3372 \\
\midrule
Peso bruto carrete (kg) & 3588 & 3582 & 3589 & 3672 \\
\midrule
Tipo de carrete & $75"\times38"\times30"$ & $75"\times38"\times30"$ & $75"\times38"\times30"$ & $75"\times38"\times30"$ \\
\midrule
Marca & CABELUM & CABELUM & CABELUM & CABELUM \\
\bottomrule
\end{tabular}
\end{center}

\vspace{0.3cm}
\textbf{Nota sobre ampacidad:} Calculada bajo condiciones estándar de CABELUM
(temperatura ambiente 25$^\circ$C, temperatura del conductor 75$^\circ$C,
viento de 0.6 m/s y radiación solar de 1000 W/m$^2$).
\end{landscape}

\section{Desglose por Norma y Tipo de Conductor}

\begin{landscape}
\subsection{Conductor AAC (All Aluminum Conductor)}

Comparativa del calibre asignado frente a los calibres más cercanos de las
normas europeas:
\begin{center}
\begin{tabular}{p{4cm}c|c|c}
\toprule
\textbf{Parámetro} & \textbf{ASTM B-231} & \textbf{BRITISH (BS 215)} & \textbf{DIN 48201-5} \\
& (954 kcmil Goldenrod) & (Tarantula) & (500 mm$^2$) \\
\midrule
Calibre / Designación & 954 kcmil & 750 mm$^2$ nominal & 500 mm$^2$ nominal \\
\midrule
Sección real (mm$^2$) & 483.4 & 794.87 & 499.83 \\
\midrule
Diámetro (mm) & 28.62 & 36.61 & 29.10 \\
\midrule
Peso (kg/km) & 1331.0 & 2192 & 1379 \\
\midrule
Carga de rotura (kgf) & 7648 & 12243 & 7615 \\
\midrule
Resistencia DC 20$^\circ$C ($\Omega$/km) & 0.05964 & 0.03627 & 0.05782 \\
\midrule
Ampacidad (A) & 985 & 899 & 689 \\
\midrule
Long. máx. despacho (m) & 2470 & 1510 & 1120 \\
\midrule
Peso neto carrete (kg) & 3288 & 3310 & 1544 \\
\midrule
Peso bruto carrete (kg) & 3588 & 3610 & 1689 \\
\midrule
Marca & CABELUM & CABELUM & CABELUM \\
\bottomrule
\end{tabular}
\end{center}
\end{landscape}

\begin{landscape}
\subsection{Conductor AAAC (All Aluminum Alloy Conductor)}

Comparativa utilizando el calibre Greeley (ASTM) contra las mayores
capacidades de fabricación de las normas europeas:
\begin{center}
\begin{tabular}{p{4cm}c|c|c}
\toprule
\textbf{Parámetro} & \textbf{ASTM B-399} & \textbf{BRITISH (BS 3242)} & \textbf{DIN 48201-6} \\
& (927.2 kcmil Greeley) & (Araucaria) & (500 mm$^2$) \\
\midrule
Calibre / Designación & 927.2 kcmil & 700 mm$^2$ nominal & 500 mm$^2$ nominal \\
\midrule
Sección real (mm$^2$) & 470.0 & 821.15 & 499.83 \\
\midrule
Diámetro (mm) & 28.14 & 37.26 & 29.10 \\
\midrule
Peso (kg/km) & 1289.0 & 2266 & 1379 \\
\midrule
Carga de rotura (kgf) & 13767 & 23465 & 14236 \\
\midrule
Resistencia DC 20$^\circ$C ($\Omega$/km) & 0.07133 & 0.04046 & 0.0671 \\
\midrule
Ampacidad (A) & 905 & --- & --- \\
\midrule
Long. máx. despacho (m) & 2550 & 1450 & 2390 \\
\midrule
Peso neto carrete (kg) & 3282 & 3286 & 3296 \\
\midrule
Peso bruto carrete (kg) & 3582 & 3586 & 3596 \\
\midrule
Marca & CABELUM & CABELUM & CABELUM \\
\bottomrule
\end{tabular}
\end{center}
\end{landscape}

\begin{landscape}
\subsection{Conductor ACAR (Aluminum Conductor Alloy Reinforced)}

La norma BRITISH y la norma DIN \textbf{no contemplan la fabricación de la
tipología ACAR} en el catálogo de CABELUM. La planta produce esta línea de
alta potencia exclusivamente bajo el estándar norteamericano.
\begin{center}
\begin{tabular}{p{4cm}c|c|c}
\toprule
\textbf{Parámetro} & \textbf{ASTM B-524} & \textbf{BRITISH} & \textbf{DIN} \\
& (950.0 MCM 18/19) & & \\
\midrule
Calibre / Designación & 950.0 MCM (18 Al / 19 Aleac.) & No disponible & No disponible \\
\midrule
Sección real (mm$^2$) & 481.0 & --- & --- \\
\midrule
Diámetro (mm) & 28.48 & --- & --- \\
\midrule
Peso (kg/km) & 1321.0 & --- & --- \\
\midrule
Carga de rotura (kgf) & 10810 & --- & --- \\
\midrule
Resistencia DC 20$^\circ$C ($\Omega$/km) & 0.06490 & --- & --- \\
\midrule
Ampacidad (A) & 1044 & --- & --- \\
\midrule
Long. máx. despacho (m) & 2490 & --- & --- \\
\midrule
Peso neto carrete (kg) & 3289 & --- & --- \\
\midrule
Peso bruto carrete (kg) & 3589 & --- & --- \\
\midrule
Marca & CABELUM & --- & --- \\
\bottomrule
\end{tabular}
\end{center}
\end{landscape}

\begin{landscape}
\subsection{Conductor ACSR (Aluminum Conductor Steel Reinforced)}

Comparativa del calibre Rail (954 kcmil) contra sus equivalentes en
dimensiones físicas europeas:
\begin{center}
\begin{tabular}{p{4cm}c|c|c}
\toprule
\textbf{Parámetro} & \textbf{ASTM B-232} & \textbf{BRITISH (BS 215 P.2)} & \textbf{DIN 48204} \\
& (954 kcmil Rail) & (Moose) & (550/70) \\
\midrule
Calibre / Designación & 954.0 kcmil & 500 mm$^2$ nominal & 550/70 mm$^2$ nominal \\
\midrule
Sección real (mm$^2$) & 483.85 & 597.00 & 620.90 \\
& (Al: 450.26 / Ac: 33.59) & & \\
\midrule
Diámetro (mm) & 29.61 & 31.77 & 32.40 \\
\midrule
Peso (kg/km) & 1598 & 1996 & 2085 \\
& (Al: 1337 / Ac: 261) & & \\
\midrule
Carga de rotura (kgf) & 11749 & 16417 & 17073 \\
\midrule
Resistencia DC 20$^\circ$C ($\Omega$/km) & 0.05988 & 0.05470 & 0.0526 \\
\midrule
Ampacidad (A) & 990 & 763 & 734 \\
\midrule
Long. máx. despacho (m) & 2110 & 1690 & 1620 \\
\midrule
Peso neto carrete (kg) & 3372 & 3377 & 3378 \\
\midrule
Peso bruto carrete (kg) & 3672 & 3677 & 3678 \\
\midrule
Marca & CABELUM & CABELUM & CABELUM \\
\bottomrule
\end{tabular}
\end{center}
\end{landscape}

\section{Análisis Comparativo}

\subsection{Por tipo de conductor}

Se observa que el conductor \textbf{AAC Goldenrod} (954 kcmil) presenta la
menor resistencia eléctrica (0.05964 $\Omega$/km), lo que lo hace el más
eficiente desde el punto de vista eléctrico, pero su carga de rotura es la
más baja (7648 kgf), limitando su uso a aplicaciones de corta distancia.

El \textbf{AAAC Greeley} (927.2 kcmil) ofrece la carga de rotura más alta del
grupo (13767 kgf) gracias a la aleación 6201-T81, aunque su resistencia
eléctrica es la mayor (0.07133 $\Omega$/km) y su ampacidad la más baja (905 A).

El \textbf{ACAR 950 kcmil 18/19} representa el mejor equilibrio para transmisión
de alta potencia. Consigue la ampacidad más alta (1044 A) manteniendo una
excelente carga de rotura (10810 kgf), siendo la opción más recomendable.

El \textbf{ACSR Rail} (954 kcmil), con su núcleo de acero de 7 hilos, ofrece
una carga de rotura de 11749 kgf y una ampacidad de 990 A, con el mayor peso
(1598 kg/km) debido al aporte del acero. Su resistencia DC es de 0.05988
$\Omega$/km, ligeramente superior a la del AAC.

\subsection{Por norma de fabricación}

Las normas \textbf{DIN} y \textbf{BRITISH} (BS) no contemplan la denominación
MCM para calibres grandes, ya que trabajan con secciones en mm$^2$. La
comparativa real se efectúa bajo el estándar \textbf{ASTM}, que es el que rige
la producción de CABELUM para estos calibres.

En las tablas de la Sección~5 se aprecia que, para cada tipo de conductor, el
calibre ASTM asignado (950 MCM) se compara contra los calibres físicamente más
cercanos disponibles en normas BRITISH y DIN:

\begin{itemize}
    \item \textbf{AAC:} 954 kcmil (ASTM) vs Tarantula 750 mm$^2$ (BS) vs
        500 mm$^2$ (DIN). La sección más cercana a 483.4 mm$^2$ corresponde
        a la norma DIN 500 mm$^2$.
    \item \textbf{AAAC:} 927.2 kcmil (ASTM) vs Araucaria 700 mm$^2$ (BS) vs
        500 mm$^2$ (DIN). El DIN 500 mm$^2$ es nuevamente el más próximo.
    \item \textbf{ACAR:} Solo disponible bajo norma ASTM. Las normas BS y DIN
        no contemplan esta tipología de conductor en el catálogo.
    \item \textbf{ACSR:} 954 kcmil (ASTM) vs Moose 500 mm$^2$ (BS) vs
        550/70 mm$^2$ (DIN). La ampacidad del ASTM (990 A) supera
        significativamente a la BS (763 A) y DIN (734 A).
\end{itemize}

\subsection{Conductor asignado: Emerson 950 MCM}

El conductor de tipo \textbf{ACAR 950 kcmil 18/19} es el recomendado para
aplicaciones de transmisión de alta potencia, ofreciendo el mejor balance
entre capacidad de corriente (1044 A) y resistencia mecánica (10810 kgf),
con una sección de 481 mm$^2$ y un diámetro de 28.48 mm, fabricado por
CVG CABELUM bajo norma ASTM B-524.

\section{Contacto con Proveedor --- CVG CABELUM}

\subsection{Datos de Contacto}

\begin{center}
\begin{tabular}{ll}
\toprule
\textbf{Dependencia} & \textbf{Datos} \\
\midrule
Oficina Comercial Caracas & Av. La Estancia, Edf. General, Piso 4, Chuao \\
& Teléfonos: 0414-103.16.55 / 0416-834.25.29 \\
\midrule
Planta Ciudad Bolívar & Vía Perimetral, Nueva Zona Industrial \\
& Teléfonos: 0285-444.35.22 / 0416-685.12.18 \\
\midrule
Correos electrónicos & comercializacion@cabelum.com.ve \\
& contacto@cabelum.com.ve \\
\bottomrule
\end{tabular}
\end{center}

\subsection{Condiciones Comerciales}

\begin{center}
\begin{tabular}{ll}
\toprule
\textbf{Aspecto} & \textbf{Información} \\
\midrule
Precio & No es fijo. Se cotiza bajo Fórmula de Reajuste \\
& de Precios indexada al valor del metal en la \\
& bolsa de la LME (London Metal Exchange). \\
\midrule
Condiciones de pago & Anticipo del 50\% al 70\% para proceder al \\
& trefilado y cableado. Saldo restante contra \\
& informes de inspección de calidad en fábrica. \\
\midrule
Condiciones de entrega & FOB Planta (Incoterms) en Ciudad Bolívar u \\
& O.A.D (Puesto en Obra) según se acuerde. \\
\midrule
Plazo de entrega & 45 a 60 días, dependiendo de la cola de \\
& producción y disponibilidad de alambrón. \\
\bottomrule
\end{tabular}
\end{center}

\section{Conclusiones Técnicas}

\begin{enumerate}
    \item \textbf{Comportamiento eléctrico frente al mecánico:} El conductor de
        tipo \textbf{ACAR 950 kcmil 18/19} es el que ofrece el mejor equilibrio
        para transmisión de alta potencia. Consigue la ampacidad más alta del
        grupo (1044 A) manteniendo una excelente carga de rotura (10810 kgf)
        gracias a sus hilos intercalados de aleación 6201.

    \item \textbf{Limitaciones del catálogo:} Los conductores \textbf{AACSR} no
        forman parte de la línea comercial regular de CABELUM. Asimismo, las
        normas \textbf{DIN} y \textbf{BRITISH} no contemplan la denominación
        ``MCM'' para calibres grandes, trabajando exclusivamente con secciones
        en mm$^2$. La totalidad de los conductores de gran sección se fabrican
        bajo normas \textbf{ASTM}.

    \item \textbf{Proveedor nacional:} CVG CABELUM es la única empresa
        venezolana que fabrica conductores de alta potencia, operando bajo
        estándares ASTM con capacidad de producción en su planta de Ciudad
        Bolívar. Los cuatro conductores evaluados (AAC Goldenrod, AAAC Greeley,
        ACAR 18/19 y ACSR Rail) son fabricados por CABELUM.

    \item \textbf{Condiciones comerciales:} Al tratarse de una empresa del
        Estado venezolano, los precios se indexan al LME y se requiere un
        anticipo del 50\%-70\% para órdenes de compra institucionales, con
        plazos de entrega de 45 a 60 días. El empaque estándar para estos
        calibres corresponde a carretes de 75" × 38" × 30", permitiendo
        6 carretes por contenedor de 20' u 8 por gandola.
\end{enumerate}

\section{Referencias}

\begin{enumerate}
    \item ASTM B-231 --- Especificación estándar para conductores AAC de aluminio 1350
    \item ASTM B-399 --- Especificación estándar para conductores AAAC de aleación 6201-T81
    \item ASTM B-524 --- Especificación estándar para conductores ACAR
    \item ASTM B-232 --- Especificación estándar para conductores ACSR
    \item CVG CABELUM --- Catálogo técnico de conductores desnudos, Planta Ciudad Bolívar
    \item British Standards BS 215 / BS 3242 --- Conductores para líneas de transmisión
    \item DIN 48201 --- Norma alemana para conductores eléctricos
\end{enumerate}

\end{document}
